\clearpage
\section{Example Outputs}

The following examples summarize the validated behaviour on deterministic test
datasets.

\subsection{Empty Directory}

For an empty directory, the report contains zero files and zero counts in all
bands.

\begin{verbatim}
Total files: 0
Bands: [0, 0, 0, 0, 0]
\end{verbatim}

\subsection{Flat Directory With Boundary Sizes}

With \texttt{MaxFS = 100} and \texttt{NB = 4}, the tested flat dataset produced
the following distribution:

\begin{verbatim}
Total files: 8
Bands: [2, 1, 1, 2, 2]
\end{verbatim}

This dataset includes files sized 0, 10, 25, 50, 75, 100, 101, and 250 bytes.

\subsection{Nested Hierarchy}

The nested test fixture confirmed that recursive traversal counts files in
subdirectories correctly and produces identical results in all three
implementations.

\begin{verbatim}
Total files: 5
Bands: [1, 1, 1, 1, 1]
\end{verbatim}

\subsection{CLI Shape}

A typical CLI run prints the selected configuration first, then rewrites the
progress line while reports are produced, and finally prints a compact table:

\begin{verbatim}
Starting CLI scan using paradigm: VT
Directory: /path/to/root
Max Size Threshold: 10.0 MiB (10485760 bytes)
Number of bands: 5
...
FINAL FILE SIZE DISTRIBUTION REPORT
Size Range Band                | File Count
\end{verbatim}

The same report object drives the GUI table, so the presentation changes but
the computed data does not.
